OSC based EOG sensor driver


This module receives data from an open hardware electro-oculography (EOG) sensor. It connects to the sensor and configures the sampling rate. Based on the data, it supports three stages: in the first stage, the data can be recorded using the data logging module. The second stage provides a post-filter (see variables starting with ``pf\_'') consisting of a low-pass filter and a maximum and minimum tracker to detect eye blinks independently of drift. The third stage is the animation of the blinks, which can be either ``none'', ``transmitted'' or ``randomised''.

\input{modtabtex/tabosceog.tex}

\input{oscdoc_tascarmod_osceog.tex}

In animation mode the duration of each blink $\tau_{d}$ is drawn from a log-normal distribution defined by the variables \attr{pf\_anim\_blink\_duration\_mu} and \attr{pf\_anim\_blink\_duration\_sigma}. The duration is limited to the interval $[0,20\,\textrm{s}]$ seconds The time between two blinks is taken from a normal distribution of blink frequency (\attr{pf\_anim\_blink\_freq\_mu} and \attr{pf\_anim\_blink\_freq\_sigma}) and transformed into a duration. The period time is limited to the interval $[\tau_{d},30\,\textrm{s}]$ to avoid excessively long blink gaps. The rectangular eye blink data is filtered with an attack-release filter with the time constants \attr{pf\_anim\_random\_tau\_attack} and \attr{pf\_anim\_random\_tau\_release}.

\subsection*{ESP8266 EOG sensor communication protocol}

This section describes the Open Sound Control (OSC) communication protocol for the Electrooculography (EOG) sensor based on the ESP8266 platform and ADS1115 ADCs. The sensor operates as a WiFi Access Point by default but can be configured to join an existing network. It streams biopotential data via UDP.


\subsection*{Network Operation Modes}

{\bf Default Access Point Mode:} By default, the device creates a WiFi Access Point.

\begin{tabular}{|ll|}
\hline
{SSID:}       & \texttt{head-<MAC\_LAST\_4\_HEX>} (e.g., \texttt{head-a1b2c3d4}) \\
{Password:}   & \texttt{headtracker}                                             \\
{IP Address:} & \texttt{192.168.100.1}                                           \\
{UDP Port:}   & \texttt{9999}                                                    \\
\hline
\end{tabular}

{\bf Station Mode (Client):}
The device can connect to an external WiFi router using the \texttt{/wlan/connect} command. Once connected, it will send data to the specified remote IP and port.

\subsection*{OSC Communication Protocol}

The device listens for incoming OSC commands on UDP port 9999. The following sections detail the supported commands.



{\bf Command: /wlan/connect}\\
Configures the WiFi connection. Can be used to switch between Access Point mode and Station mode, or to configure a remote data destination.

{\small
\begin{tabular}{|l|l|l|}
\hline
\textbf{\#} & \textbf{Type} & \textbf{Description} \\ \hline
0 & String & SSID of the network to join (Empty string to reset to AP mode) \\ \hline
1 & String & Password for the network \\ \hline
2 & String & Remote IP address (destination for EOG data) \\ \hline
3 & Integer & Remote Port (destination for EOG data) \\ \hline
\end{tabular}
}

Behavior: If called with 0 arguments: Resets to default Access Point mode. If called with 4 arguments: Connects to the specified WiFi network and configures the destination for data streaming. Upon successful connection in Station mode, the outgoing OSC address path is automatically updated to \texttt{/<MAC\_LAST\_4\_HEX>/eog}.


{\bf Command: /eog/connect}\\
Configures the destination for EOG data without changing WiFi settings. This is typically used when the device is already in Station mode or to redirect data to a different listener on the current network.

{\small
\begin{tabular}{|l|l|l|}
\hline
\textbf{\#} & \textbf{Type} & \textbf{Description} \\ \hline
0 & Integer & Remote Port \\ \hline
1 & String & OSC Address Path (e.g., \texttt{/eog}) \\ \hline
\end{tabular}
}

\textbf{Note:} The IP address used is the IP of the last device that sent a command to the sensor (source IP of the incoming UDP packet).

{\bf Command: /eog/disconnect}\\
Stops the streaming of EOG data. Arguments: None. Effect: Sets the
remote port to 0, effectively disabling data transmission.

{\bf Command: /eog/srate}\\
Sets the sampling rate of the ADCs.

{\small
\begin{tabular}{|l|l|l|}
\hline
\textbf{\#} & \textbf{Type} & \textbf{Description} \\ \hline
0 & Int/Float/Double & Desired sampling rate (Hz) \\ \hline
\end{tabular}
}

{\bf Supported Ranges and ADC Settings:}
The command maps the requested rate to the nearest supported ADS1115 data rate.

{\small
\begin{tabular}{|l|l|}
\hline
\textbf{Input Range (Hz)} & \textbf{Actual Rate (SPS)} \\ \hline
$< 16$ & 8 \\ \hline
$16 - 31$ & 16 \\ \hline
$32 - 63$ & 32 \\ \hline
$64 - 127$ & 64 \\ \hline
$128 - 249$ & 128 \\ \hline
$250 - 474$ & 250 \\ \hline
$475 - 859$ & 475 \\ \hline
$\ge 860$ & 860 \\ \hline
\end{tabular}
}

{\bf Command: /eog/urange}\\
Sets the input voltage range of the ADC.

{\small
\begin{tabular}{|l|l|l|}
\hline
\textbf{\#} & \textbf{Type} & \textbf{Description} \\ \hline
0 & Float & Maximum voltage range in Volts \\ \hline
\end{tabular}
}

{\bf Supported Ranges:}
The command maps the input value to the closest ADS1115 gain setting.

{\small
\begin{tabular}{|l|l|}
\hline
\textbf{Input Range (V)} & \textbf{ADC Setting} \\ \hline
$< 0.257$ & $\pm 0.256V$ \\ \hline
$0.257 - 0.512$ & $\pm 0.512V$ \\ \hline
$0.513 - 1.024$ & $\pm 1.024V$ \\ \hline
\end{tabular}
}

{\bf Data Streaming:} Once connected via \texttt{/eog/connect}
or \texttt{/wlan/connect}, the device will stream sensor data.  The
OSC Address is configurable via \texttt{/eog/connect}, and defaults
to \texttt{/<MAC\_LAST\_4\_HEX>/eog}. The message contains 3 floats:

{\small
\begin{tabular}{|l|l|l|}
\hline
\textbf{\#} & \textbf{Type} & \textbf{Description} \\ \hline
0 & Float & Timestamp (seconds) \\ \hline
1 & Float & Voltage Channel 1 (Volts) \\ \hline
2 & Float & Voltage Channel 2 (Volts) \\ \hline
\end{tabular}
}

{\bf Timing:}
The device attempts to send data at the interval defined by \texttt{sample\_period} ($1 / \text{srate}$). The timestamp is derived from \texttt{micros()}.
